\chapter{K-means clustering and quantization}

\section{K-means formulation}

Clustering is a fundamental problem in unsupervised learning, where the goal
is to partition a given dataset into meaningful groups. The intuitive image is
given in Figure~\ref{fig:ch8-three-class} where initially only the left image
(un-colored points) is given; the clustering algorithm has to assign each point
to a cluster which is here figured by a different color. Note that here the
overall goal is much less specified than in the case of classifications: there
are no explicit indication of \textbf{how to partition} the points (no training
set), no \textbf{ground truth to validate the results} and sometimes the number
of partitions is not known. The task is not always as the Figure~\ref{fig:ch8-three-class}
may suggest, consult for instance Figure 1.2 where one can see that for some
examples such as those in first two lines intuitive behavior is \textbf{not}
recovered by simple clustering algorithms.

One of the most widely used clustering algorithms is K-Means, which seeks to
minimize intra-cluster variance.

\begin{definition}[K-Means Clustering]
Given a dataset $\cX = \{x_1, x_2, \ldots, x_n\} \subset \R^d$ and an integer $K$,
a K-Means clustering of $\cX$ is a partition $C = (C_1, C_2, \ldots, C_K)$ of
$\cX$\footnote{Recall that $C = (C_1, C_2, \ldots, C_K)$ is said to be a
partition of $\cX$ if $C_i \cap C_j = \emptyset$ for all $i \le j \le K$ and
$\cup_{k=1}^{K} C_k = \cX$.} into $K$ disjoint clusters $C_1, C_2, \ldots, C_K$
such that the within-cluster sum of squared distances to the mean is minimized:
\begin{equation}\label{eq:ch8-kmeans-obj}
\min_{C_1, \ldots, C_K} \sum_{k=1}^{K} \sum_{x \in C_k} \norm{x - \mu_k}^2,
\end{equation}
where $\mu_k = \frac{1}{|C_k|} \sum_{x \in C_k} x$ is the centroid of the cluster
$C_k$. Here $\norm{\cdot}$ is the usual Euclidian norm in $\R^d$.
\end{definition}

\begin{figure}[h]\centering
% FIGURE PLACEHOLDER (uncomment & replace with \includegraphics when image is available):
% \fbox{\parbox{0.7\linewidth}{\centering\textit{[Figure from PDF page 116:
% two scatter plots showing a three-class clustering example. Left: unlabeled
% blue points in three groups. Right: points colored by cluster (blue, red, green)
% with dashed Voronoi boundary lines.]}}}
\caption{Three class clustering example together with its Voronoi partition.
K-means input: left figure; K-means output: the right figure.}
\label{fig:ch8-three-class}
\end{figure}

The explanation for the terminology and for choosing $\mu_k = \frac{1}{|C_k|}
\sum_{x \in C_k} x$ is given in the following result:

\begin{lemma}
Let $\mathcal{P} = \{p_1, \ldots, p_n\} \subset \R^d$. The problem
$\min_{\mu \in \R^d} \sum_{p \in \mathcal{P}} \norm{p - \mu}^2$ has an unique
minimum $\mu^\star = \frac{1}{n} \sum_{p \in \mathcal{P}} p$ and:
\begin{equation}\label{eq:ch8-centroid}
\min_{\mu \in \R^d} \sum_{p \in \mathcal{P}} \norm{p - \mu}^2
= \sum_{p \in \mathcal{P}} \norm{p - \mu^\star}^2
= \frac{1}{2n} \sum_{p, p' \in \mathcal{P}} \norm{p - p'}^2
= n \Var[\mathcal{P}].
\end{equation}
\end{lemma}

\begin{proof}
The proof is obtained by writing the critical point equations or using
bias-variance decomposition:
\begin{equation}\label{eq:ch8-bias-var}
\frac{1}{n} \sum_{p \in \mathcal{P}} \norm{p - \mu}^2
= \Var[\mathcal{P}] + \norm{\mu - \mu^\star}^2.
\end{equation}
\end{proof}

The intuitive vision of a clustering problem is that well defined clusters
exist, that is, some partition in clusters $C_i$ exits such that distances
between members of same cluster are small and distances between elements from
different clusters are large. The result below ensures that the clustering
procedure will recover intuitive results.

\begin{lemma}
Assume that the dataset $\cX = \{x_1, x_2, \ldots, x_n\}$ has a partition
$C_1^\star, \ldots, C_K^\star$ with $C_i^\star \neq \emptyset$ for $i \le K$
and there exists two constants $d_m$ and $d_M$ such that:
\begin{align}
\text{small intra cluster distances: } & \forall k \le K, \forall x, y \in C_k^\star :
\norm{x - y} \le d_m \label{eq:ch8-intra} \\
\text{large inter cluster distances: } & \forall k \neq \ell \le K,\ \forall x \in C_k^\star, z \in C_\ell^\star :
\norm{x - z} \ge d_M. \label{eq:ch8-inter}
\end{align}
If $d_M > d_m \sqrt{2n}$ the solution to the K-means problem is exactly the
partition $C_1^\star, \ldots, C_K^\star$ (up to a relabeling of the classes).
\end{lemma}

\begin{proof}
Note that if $x \in C_\ell^\star$ and $C_\ell^\star$ has mean $\mu_\ell^\star$
then:
\begin{equation}\label{eq:ch8-mean-bound}
\norm{x - \mu_\ell^\star}
= \frac{1}{|C_\ell^\star|} \norm[\bigg]{\sum_{y \in C_\ell^\star} (x - y)}
\le \frac{1}{|C_\ell^\star|} \sum_{y \in C_\ell^\star} \norm{x - y}
\le \frac{|C_\ell^\star|}{|C_\ell^\star|} d_m = d_m.
\end{equation}
Thus the value of the partition $C_1^\star, \ldots, C_K^\star$ is at most
$n d_m^2$; on the other hand, if we have a different partition $C_1, \ldots, C_K$
there will be some $\ell$ and $x \neq y, x, y \in C_\ell$ but such that $x, y$
are not in the same set of the partition $C^\star$. Denote $\mu_\ell$ the
centroid of $C_\ell$; the value of the partition $C$ is at least
$\norm{x - \mu_\ell}^2 + \norm{y - \mu_\ell}^2 \ge \frac{1}{2} \norm{x - y}^2
\ge \frac{d_M^2}{2}$. Thus, the value of the partition $C$ is larger than
$d_M^2 / 2$ which, by hypothesis, is strictly larger than $n d_m^2$ which is
larger than the value of the partition $C^\star$; this proves that $C^\star$
is the optimal partition.
\end{proof}

\begin{lemma}[Strict improvement with increasing $K$, different points]
Let $\cX = \{x_1, \ldots, x_n\} \subset \R^d$ be a finite collection of points.
Denote, for $\mu = (\mu_1, \ldots, \mu_K)$:
\begin{equation}\label{eq:ch8-Jmu}
J(\mu) = \frac{1}{n} \sum_{x \in \cX} \min_{1 \le k \le K} \norm{x - \mu_k}^2,
\end{equation}
and let $J_K^\star$ denote the minimal K-means cost, i.e.,
\begin{equation}\label{eq:ch8-JKstar}
J_K^\star = \min_{\mu \in (\R^d)^K} J(\mu).
\end{equation}
Then
\begin{enumerate}
    \item $J_{K+1}^\star \le J_K^\star$, with equality if and only if $\cX$
    has exactly $K$ distinct points.
    \item In particular, if $\cX$ has at least $K$ distinct points then the
    centroids $\mu_1, \ldots, \mu_K$ are all distinct.
\end{enumerate}
\end{lemma}

\begin{proof}
First, note that when we move from $K$ clusters to $K+1$ clusters, we have
more flexibility in the placement of the centers. Specifically, if we keep
$\mu_1, \ldots, \mu_K$ as before and place the new center $\mu_{K+1}$ arbitrarily
the cost for the new solution will not increase. This gives us the inequality
$J_{K+1}^\star \le J_K^\star$. Next, if the data points are not perfectly
partitioned into $K$ clusters with zero variance there exists a cluster with
two distinct points $x_a \neq x_b$; we can split this cluster into two new
clusters, one containing $x_a$ with centroid in $x_a$ and the other the rest
of the points with previous centroid. By doing so, we will strictly reduce the
within-cluster variance for the points that were originally assigned to the
cluster being split. Hence, the total cost decreases strictly, and we obtain
the strict inequality $J_{K+1}^\star < J_K^\star$.
Finally, using the proposition for $K-1$ and $K$ clusters we obtain that the
points have to be different.
\end{proof}

\subsection{Optimal Transportation framework}

Assume that $x_i$ are drawn from some distribution $\nu$ with support on space
$\cX$; we can formulate the K-means clustering problem as an optimal
transportation problem \cite{turinici25} between the measure $\nu$ and a
measure with a discrete support of at most $k$ points. One can prove that this
is equivalent to finding a measurable function $\mathcal{T} : \cX \to \cX$
with $\mathrm{Im}(\mathcal{T}) \subset \cX$ a discrete set of cardinality at
most $K$ that minimizes the \emph{total transportation cost}:
\begin{equation}\label{eq:ch8-OT-cost}
\E_{X \sim \nu} \norm{X - \mathcal{T}(X)}^2
= \int_{\cX} \norm{x - \mathcal{T}(x)}^2 \nu(dx).
\end{equation}
Here $(x, y) \mapsto \norm{x - y}^2$ is called the transportation cost between
locations $x$ and $y$. In this formulation, the map $\mathcal{T}$ assigns to
each point $x$ its corresponding centroid. Since image of $\mathcal{T}$ is
discrete the map $\mathcal{T}$ cannot be continuous.

\section{K-means algorithm}

To solve problem~\eqref{eq:ch8-kmeans-obj} many approaches can be proposed, but
historically an iterative refinement technique has been employed. It became so
popular that it was often called just ``the K-means algorithm''; Lloyd used it
at Bell Labs as early as 1957 but formalized it only in 1982~\cite{turinici12};
independent proposals were also developed during that time. The procedure is
described in Algorithm~\ref{alg:ch8-kmeans} and works by iteratively assigning
points to centroids and then updating the centroids to be the center of mass
of the assigned points. Then a new iteration is started.

\begin{algorithm}
\caption{K-Means Algorithm}\label{alg:ch8-kmeans}
\begin{algorithmic}[1]
\State Initialize $K$ cluster centroids $\mu_1, \ldots, \mu_K$ randomly.
Define a tie break protocol.
\Repeat
    \State Assign each point $x \in \cX$ to the closest centroid, apply the
    tie break protocol if needed; construct a partition $C_1, \ldots, C_K$ of
    $\cX$ such that:
    \begin{equation}\label{eq:ch8-assign}
    \forall k \le K,\ \forall x \in C_k : \norm{x - \mu_k}
    = \min_{\ell \le K} \norm{x - \mu_\ell}.
    \end{equation}
    \State Update each centroid:
    \begin{equation}\label{eq:ch8-update}
    \forall k \le K :\ \mu_k = \frac{1}{|C_k|} \sum_{x \in C_k} x.
    \end{equation}
\Until{Centroids do not change.}
\end{algorithmic}
\end{algorithm}

The algorithm needs a tie-break protocol, that is a way to choose, for some
point $x$ the cluster to which it belongs in case there are several possible
candidates; this arises when there exist two centroids $\mu_a$ and $\mu_b$
with: $\norm{x - \mu_a} = \norm{x - \mu_b} \le \norm{x - \mu_k}$ for all
$k \le K$. Possible protocols are random choosing or some deterministic rule;
some examples are given in Table~\ref{tab:ch8-tiebreak}. For instance the
lowest index rule is naturally obtained if using some `argmin' in Python, as
`argmin' will return first index in case of ties.

\begin{table}[h]\centering
\begin{tabular}{@{}ll@{}}
\toprule
Method & Description \\
\midrule
Arbitrary Assignment & Assign the point randomly to one of the tied clusters. \\
Lowest Index Rule & Assign the point to the cluster with the lowest index. \\
Previous Assignment & \begin{tabular}[t]{@{}l@{}}
If the point was previously assigned to one of the \\
tied clusters, it remains in that cluster, otherwise \\
random.
\end{tabular} \\
Density-Based & \begin{tabular}[t]{@{}l@{}}
Assign to the cluster with higher density (i.e., \\
having more points).
\end{tabular} \\
Weighted Probability & Assign probabilistically based on cluster sizes. \\
Perturbation & Add slight noise to break ties naturally. \\
\bottomrule
\end{tabular}
\caption{Some common tie-breaking methods in K-Means clustering.}
\label{tab:ch8-tiebreak}
\end{table}

\begin{theorem}[Convergence]
Assume a deterministic tie-breaking rule such as the lowest index rule. Then
the K-Means algorithm terminates in a finite number of iterations.
\end{theorem}

\begin{proof}
Assume there are $n$ points and define the `inertia' function:
\begin{equation}\label{eq:ch8-inertia}
J(\mu, C) = \frac{1}{n} \sum_{k=1}^{K} \sum_{x \in C_k} \norm{x - \mu_k}^2.
\end{equation}
Here $\mu = (\mu_1, \ldots, \mu_K)$ is the $K$-tuple of all centroids in the
index order and $C = (C_1, \ldots, C_K)$ is the partition written as a $K$-tuple.
Note that the inertia decreases at every step of the algorithm: this is obvious
for the assignment phase and for the updating phrase this results from the
fact that the mean $\bar{y}$ of vectors $y_1, \ldots, y_J$ is precisely the
unique minimum of $z \mapsto \sum_{j=1}^{J} \norm{y_j - z}^2$, cf.
Lemma~\ref{eq:ch8-centroid}. Moreover, due to the deterministic tie-breaking
rule, if at some point the centroid assignment $C$ remains the same from one
iteration to the next then they will never change again and convergence is
obtained. Also note that the decrease in inertia is strict, unless convergence
occurs. Since $C$ lives in a finite set there are bound to be repetitions
during the iterations of $C$ which indicate convergence.
\end{proof}

\begin{advancedtopics}
See \cite{turinici17, turinici3, turinici8} for more in-depth considerations
on the convergence of the K-means algorithm.
\end{advancedtopics}

\begin{warning}
The above result guarantees that the algorithm stops not that it finds the
global optimum. See Figure~\ref{fig:ch8-init} for an example where poor
initialization leads to suboptimal results. In practice such cases are not
exceptional that local (not global) optimums are reached. Note that this
suboptimal solution is locally stable around this point: take for instance 10
points in each vertex of the rectangle and assign some of them to the other
class, the solution will come back to the suboptimal choice. A common cure is
to run the algorithm several times from several initial points in the hope that
one may find the true global solution. Other initialization choices include
Forgy method which randomly chooses $k$ observations from the dataset and uses
these as the initial means. Modern variants use K-Means++ initialization which
improves convergence (it is based on a sampling depending on the distance).
\end{warning}

\begin{example}
The prototypical situation is when the dataset is a mixture with disjoint
support. When the supports are far enough, the centroids will be the means of
each variable in the mixture, see Lemma~\ref{eq:ch8-intra}.
\end{example}

\subsection{Further Practical Considerations}

Beyond initialization, several aspects determine the algorithm performance:
\begin{itemize}
    \item \textbf{Choice of $K$:} The number of clusters $K$ can be determined
    using the Elbow Method\footnote{In the elbow method we plot how the
    within-cluster variance decreases as $K$ increases and look for a point,
    ``elbow'' where adding more clusters doesn't improve much.} or Silhouette
    Analysis\footnote{In silhouette analysis, for each point, we compute a
    score comparing its distance to its own cluster and to the nearest other
    cluster, then average these scores over all points, and we select the
    number of clusters $K$ that maximizes this average score; this naturally
    avoids choosing a $K$ that is too large because isolated or tiny clusters
    lower the silhouette quality.}.
    \item \textbf{Computational Efficiency:} The standard algorithm has a
    worst-case exponential complexity, but is efficient in practice. Many
    variants improve this score by different means including parallelism and
    various approximations;
\end{itemize}

\begin{figure}[h]\centering
% FIGURE PLACEHOLDER (uncomment & replace with \includegraphics when image is available):
% \fbox{\parbox{0.7\linewidth}{\centering\textit{[Figure from PDF page 122:
% two side-by-side purple rectangles with four black dots at the corners.
% Left rectangle: two red centroids $\mu_1$ and $\mu_2$ placed vertically
% (top-middle and bottom-middle) indicating a sub-optimal initialization.
% Right rectangle: two red centroids $\mu_1$ and $\mu_2$ placed horizontally
% (left-middle and right-middle) indicating the optimal solution.]}}}
\caption{Importance of the initialization in K-means algorithm. We take $n=4$,
$d=2$, $k=2$; the dataset $\cX$ is the set of black points that form a
rectangle centered at the origin. \textbf{Left:} sub-optimal initialization.
The algorithm terminates after one step but the solution is not optimal.
\textbf{Right:} the optimal solution.}
\label{fig:ch8-init}
\end{figure}

\begin{itemize}
    \item \textbf{Model assumptions and behavior:} The K-means tends to produce
    clusters of similar size because minimizing the total squared distances
    induces strong penalization for large distances and thus favors compact,
    evenly distributed groups, and the Voronoi partition naturally balances
    the number of points assigned to each centroid.
\end{itemize}

\section{Other distances}

Similar questions can be asked for different minimization settings. For instance,
replacing the distance squared in the inertia with just the distance allows to
formulate the generalized Weber problem:

\begin{definition}[Multi-Weber Problem]
Given a collection of points $\cX = \{x_1, x_2, \ldots, x_n\} \subset \R^d$
and an integer $K$, the solution of the generalized Weber problem is a
$K$-tuple of points (centroids) $\mu = (\mu_1, \mu_2, \ldots, \mu_K)
\in (\R^d)^K$ that minimize the total sum of distances:
\begin{equation}\label{eq:ch8-weber}
\min_{\mu} \sum_{k=1}^{K} \sum_{x \in C_k} \norm{x - \mu_k},
\end{equation}
where each point $x \in \cX$ is assigned to the closest centroid $\mu_k$ as
in~\eqref{eq:ch8-assign}.
\end{definition}

\begin{remark}
When $K = 1$ the question is to find $\mu$ that minimizes distance to all
$x_i$. In one dimension ($d = 1$) this is nothing else than the median and for
$d > 1$ this is one of the definitions of the median, also called the
`barycenter' of the points, cf. Exercise 8.3.
\end{remark}

\begin{advancedtopics}
The $k=1$, $d=2$, $n=3$ case is called Weber-Fermat problem and has a
geometric solution:

\begin{lemma}[Fermat-Torricelli]
If the angles of the triangle formed by $x_1, x_2, x_3$ are all less than
$120^\circ$, then the solution $\mu_1^\star$ is the unique point such that
the angles subtended by the segments from $\mu^\star$ to each $x_i$ are equal
to $120^\circ$.
\end{lemma}

The problem does not have a closed-form solution for $n \ge 4$, but numerical
optimization is used instead.
\end{advancedtopics}

\section{Voronoi Diagrams}

An important related notion is the Voronoi diagram (or Voronoi partition or
Voronoi map) of a collection of points. We will explain it for the case of
the plane ($d = 2$). A Voronoi diagram is a partitioning of a plane into
regions based on the distance to a specified set of points. Given a set of
distinct points in the plane, the Voronoi diagram divides the plane into
regions such that each region corresponds to one of the points. The region
associated with a point consists of all locations that are closer to that
point than any other. See Figure~\ref{fig:ch8-three-class} for an illustration.
In other words, for a collection of points $P = \{p_1, p_2, \ldots, p_n\}
\subset \R^d$ the Voronoi diagram assigns to each point $p_i$ a region $R_i$,
called cell, with:
\begin{equation}\label{eq:ch8-voronoi}
R_i = \{q \in \R^d \mid \norm{q - p_i} < \norm{q - p_j},\ \forall j \neq i\}
\end{equation}
where $\norm{q - p_i}$ is the Euclidean distance between a point $q$ and the
point $p_i$. Voronoi diagrams are used in clustering algorithms where Voronoi
cells correspond to the clusters around the centroids.

\section{Quantization}

K-means algorithm is an example of \emph{vector quantization}. Vector
quantization is a technique that maps a distribution of high-dimensional
vectors in $\R^d$ to a finite collection of discrete representative vectors,
still in $\R^d$, called a `codebook'. This process is commonly used for data
compression, clustering, and pattern recognition.

Sometimes a `encoding' function can be provided to associate input vectors to
a codebook vector; for instance one may construct Voronoi diagram of the
codebook and map all points in a Voronoi cell to the corresponding cell
centroid.

\begin{figure}[h]\centering
% FIGURE PLACEHOLDER (uncomment & replace with \includegraphics when image is available):
% \fbox{\parbox{0.7\linewidth}{\centering\textit{[Figure from PDF page 124:
% three images side by side. Left: a smooth full-color rainbow gradient image.
% Middle: a narrow vertical strip showing four representative colors (the
% codebook). Right: the same rainbow image quantized to only those four colors,
% producing visible vertical color bands.]}}}
\caption{Example of color quantization, see text for details. The left image
is color-quantized with the four colors in the middle; the result is the
right image.}
\label{fig:ch8-color-quant}
\end{figure}

\begin{example}
We will take a color quantization example. Consider the image on the left of
Figure~\ref{fig:ch8-color-quant}. Each color can be expressed as a 3D vector.
The size of the image is $256 \times 256$ so each pixel requires 8 bits of
color so 24 bits in all. Instead of storing all this information, we can
consider replacing the color vectors by the centroids of a clustering
algorithm. Then each color is replaced by its centroid; if we take $k = 4$ of
them we obtain the four colors depicted in the middle of the figure. Storing
such colors is much less expensive (only 2 bits of information each, so a
compression factor of 12). The resulting image is drawn in the right plot. Of
course, there is some loss of information but the general idea is preserved.
Of course, we could also consider more clusters, for instance 8 or 16 colors.
\end{example}

\begin{example}
Another example is the quantization of the Brownian motion. The Brownian
motion can be considered as a probability law on the ensemble of possible
scenarios of the Brownian motion. This is called the Wiener space view. Suppose
for instance that the time horizon is $[0, T]$ that we divide into $N$
segments and only consider times instants $t_n = nh$ where $h = T/N$. Then
the Brownian motion is reduced to the collection of increments
$W_{t_{\ell+1}} - W_{t_\ell}$, for $\ell < N - 1$. Any such increment is a
$\sqrt{h} \cN(0, 1)$ distributed random variable. So finally a scenario can
be replaced by a standard $N$-dimensional gaussian (up to a factor $\sqrt{h}$).
If we sample from this law we can consider a large sample and use K-means or
similar algorithms to find representative scenarios. In general the K-means
is not the most adapted, see~\cite{turinici23} for details. An example is
given in Figure~\ref{fig:ch8-brownian}.
\end{example}

\begin{figure}[h]\centering
% FIGURE PLACEHOLDER (uncomment & replace with \includegraphics when image is available):
% \fbox{\parbox{0.7\linewidth}{\centering\textit{[Figure from PDF page 125:
% two side-by-side line plots of Brownian motion trajectories. Left: many
% overlapping multicolored Brownian sample paths over time $[0, 60]$, forming
% a cloud that spreads from $0$ to approximately $\pm 20$. Right: $K=10$
% representative quantized trajectories on $[0, 1]$ ranging roughly within
% $[-1, 1]$.]}}}
\caption{Quantization of brownian motion. Left image: some brownian scenarios.
Right image: quantization with $N = 128$ time steps and $K = 10$ codebook.}
\label{fig:ch8-brownian}
\end{figure}
